\scanpage{74}
\uline{Un théorème d'orthogonalité}

\deleted{$X$ \unsure{schéma}}
$A$ anneau noethérien $J$-adique (séparé et complet)
\[
X=\operatorname{Spec}A\qquad
Y=\operatorname{Spec}(A/J)=V(J)
\]
\deleted{$\unsure{G|C}$}
\[
U=X-Y
\]
Pour tout schéma fini étale $Z$ sur $Y$, $\widetilde Z$ est le schéma fini étale sur $X$ (unique à isom unique près) le prolongeant.
\margin{$D=$ pour faisceaux $\ell$-adiques constructibles\unsure{,...|....}}
$G,G'$ schémas \deleted{\unsure{lisses|propres}} \inserted{en groupes comm. \deleted{\unsure{séparés}} \inserted{de t.f. séparés}} \inserted{plats} sur $X$ \deleted{séparés}
\[
\begin{array}{ll}
G_U=G|U&\text{est un schéma abélien }A\\
G'_U=G'|U&\text{\rule{5cm}{.3pt}}\quad A'
\end{array}
\]
\margin{$A$ et $A'$ \deleted{\unsure{duaux l'un de l'autre}} \inserted{accouplés} i.e. on a un élém. de corr. div. sur $A\times A'$}
$\ell$ \deleted{\unsure{entier}} nb premier inv. sur $Y$, donc sur $X$.
\margin{Commencer par parler de $\unsure{G|C},\unsure{D|G'})$ avant parler de $\unsure{\deleted{G}\inserted{A}|A},G',A)$}
$S\subset G_Y=G\times_XY$ est \inserted{plat sur $Y$ et} tel que ${}_{\ell^\nu}\unsure{S|\deleted{G}\inserted{S}}$ \uline{fini} sur $Y$ pour $\unsure{\mathrm{tt}|\forall}\nu>0$
\deleted{$\unsure{S'\subset G'_Y|G'_Y|S'_Y}=G'\times_XY\quad{}_{\ell^\nu}G'$}
\[
T'\subset G'_Y\qquad\text{un \uline{tore}}
\]
\deleted{\unsure{Alors} $\ell^\nu$}
$\widetilde{T_\ell(\deleted{G}\inserted{S})}\deleted{\unsure{\subset|e}}$ est défini, et il est $\subset T_\ell(G)$
\unsure{idem|id.} $=\widetilde{T_\ell(T)}$ \unsure{défini|defini} et $\subset T_\ell(G')$
\margin{\uline{NB} les ${}_{\ell^\nu}G$ sont etales de t.f. sép. sur $S$, et leur restriction à $Y$ contient les ${}_{\ell^\nu}S$, donc ${}_{\ell^\nu}G$ contient ${}_{\ell^\nu}\widetilde S$.}
\diagram{\begin{tikzcd}[column sep=small,row sep=small]
\widetilde{T_\ell(\deleted{G}\inserted{S})}|U\subset{} & T_\ell(A)\\[-.4em]
{} & \times\\[-.4em]
\widetilde{T_\ell(T)}|U\subset{} & T_\ell(A')\arrow[d,"\varphi_\xi"',"{\text{--- \unsure{pair|voir}: }\xi\text{ induit }T_\ell(S)\unsure{'|.}}"]\\
{} & T_\ell(\mathbf G_{mU})=\widetilde{T_\ell(\mathbf G_{mY})}|U
\end{tikzcd}}
\uline{Théorème}
\[
\varphi_\xi\bigl(\widetilde{T_\ell(S)}|U,\widetilde{T_\ell(T)}|U\bigr)=0
\]
\unsure{1} \uline{Dém} Il suffit de le prouver en chaque pt $x\in U$.
Prenant un point \inserted{$y\in\overline{\{x\}}\cap Y$} \deleted{dans l'adhérence}, et dominant par un anneau de val. discret, on est ramené au cas où $X$ est un \uline{trait}, $Y$ son pt fermé, $U$ son pt gén. Il suffit d'ailleurs, dans cette réduction, de prendre pour $x$ un pt maximal de $U$ : Supposons \deleted{\unsure{qu'|qui|que}} d'abord qu'on
\margin{73}
